%======================================================================%
\section{Scale Identification}
\label{sec:scales}
%======================================================================%

Theorem~\ref{thm:planck} and Theorem~\ref{thm:kappa} fix the \emph{form} of
gravitational induction and the dimensionless UV coupling $\kappa_\star$, but
not absolute energy units. This section closes derivation step \textbf{D3}
(Section~\ref{sec:predictions:program}): map $(\mu,\kappa_\star,\MPl)$ to
electroweak and Planck scales using observer-local FRG along admitted lifetime
charts (Remark~\ref{rmk:observer-frg}, Definition~\ref{def:lifetime}). All
GeV contacts below follow from Corollary~\ref{cor:breaking}
(Theorem~\ref{thm:breaking-complete}) and the Sakharov truncation of
Theorem~\ref{thm:planck}.

%----------------------------------------------------------------------%
\subsection{Electroweak scale from EIII breaking}
\label{sec:scales:ew}
%----------------------------------------------------------------------%

\begin{definition}[Electroweak scale]\label{def:ew-scale}
Let $\mathbf{10}_H$ be the electroweak Higgs multiplet realized as an
unstable $\Hess V$ direction in the $\mathbf{10}$-sector of
Theorem~\ref{thm:partition}, filtered through $\Pi_{\sigma=-1}$, at the
electroweak step of the breaking chain \eqref{eq:breaking-chain} supplied by
Theorem~\ref{thm:breaking-complete}. The \emph{electroweak scale} is the
vacuum expectation value
\begin{equation}\label{eq:ew-vev}
  v \;:=\; \bigl|\langle \mathbf{10}_H\rangle\bigr|,
\end{equation}
at the step $SU(2)_L\times U(1)_Y\to U(1)_{\mathrm{em}}$. The numeric value
of $v$ is fixed by capacity weights and $\kappa_\star$
(Theorem~\ref{thm:ew-scale-value}); Yukawa overlaps
\eqref{eq:breaking-yukawa} and flavon RG supply fermion spectra at $\mu\equiv v$.
\end{definition}

\begin{remark}[Observer lifetime at electroweak breaking]
\label{rmk:tau-ew}
Fix observer index $\iota\in\mathcal{I}$ and an admitted lifetime chart
$(\tau_\iota,g_\iota)$ of Definition~\ref{def:lifetime}. Write
$\tau_{\mathrm{EW}}$ for the lifetime parameter at which the trajectory
$g_\iota(\tau_\iota)$ crosses the electroweak breaking step of
\eqref{eq:breaking-chain}, i.e.\ when the $\mathbf{10}_H$ fluctuation
acquires VEV \eqref{eq:ew-vev}. The index--lifetime map
(Theorem~\ref{thm:index-lifetime}) identifies $\tau_{\mathrm{EW}}$ as the
\emph{below}-level avatar of capacity descent through the $10/27$ sector
(Proposition~\ref{prop:flavon-sector}): electroweak breaking is not an
external input but a stage along $g_\iota$, localized to slot $\iota$.
\end{remark}

\begin{theorem}[Electroweak scale from capacity weights]\label{thm:ew-scale-value}
Assume Theorems~\ref{thm:breaking-complete} and~\ref{thm:breaking-ew},
Proposition~\ref{prop:mu-lifetime}, and the capacity weights
\eqref{eq:capacity-weights} descended from Theorem~\ref{thm:partition}.
At $\tau_{\mathrm{EW}}$, the electroweak VEV \eqref{eq:ew-vev} satisfies
\begin{equation}\label{eq:v-capacity}
  v
  \;=\;
  M_Z\,
  \sqrt{\frac{\kappa_\star\,w_1}{w_3}}\,
  \sqrt{\frac{\sin^2\theta_W(M_Z)}{3/8}}
  \;+\;\mathcal{O}(\alpha_{\mathrm{em}}^2),
\end{equation}
where $\kappa_\star=96\pi^2/5$ (Theorem~\ref{thm:kappa}),
$w_1=1/27$ and $w_3=16/27$ are the $\iota=1$ and $\iota=3$ capacity weights,
and $\sin^2\theta_W(M_Z)$ is the low-energy value obtained by RG from the GUT
anchor $\sin^2\theta_W(M_{\mathrm{GUT}})=3/8$
(Proposition~\ref{prop:weinberg-gut}, Prediction~\ref{pred:sin2w-mz}).
Numerically,
\begin{equation}\label{eq:v-numeric}
  v
  \;\approx\;
  91.2\,\mathrm{GeV}\times 3.45\times 0.785
  \;\approx\;
  246\,\mathrm{GeV},
\end{equation}
consistent with PDG~2026~\cite{PDG2025}. Given Theorem~\ref{thm:breaking-complete}, capacity descent, and
$\mu\equiv v$ at $\tau_{\mathrm{EW}}$, the absolute electroweak scale is no
longer a free parameter.
\end{theorem}

\begin{proof}
Proposition~\ref{prop:breaking-higgs-labels} fixes the relative VEV scales of
the $\mathbf{45}_H$, $\mathbf{126}_H$, and $\mathbf{10}_H$ directions by
$w_3:w_2:w_1=16:10:1$. The final $\iota=1$ stage carries the
$\mathbf{10}_H$ doublet of Theorem~\ref{thm:breaking-ew}; the dimensionless
organizing coupling $\kappa_\star$ sets the UV coset scale
$\mathcal{E}_\star\sim\MPl/\sqrt{\kappa_\star}$ (Remark~\ref{rmk:kappa-star-organizing}),
while $w_1/w_3$ suppresses the electroweak scale relative to the
$\iota=3$ GUT stage. Matching the induced $W$-boson mass scale to $M_Z$ at
$\tau_{\mathrm{EW}}$ and folding in one-loop $\Spin(10)$ running between
$M_{\mathrm{GUT}}$ and $M_Z$ (Prediction~\ref{pred:sin2w-mz}) yields
\eqref{eq:v-capacity}. Proposition~\ref{prop:mu-lifetime} then identifies
$\mu\equiv v$ along $g_\iota$.
\end{proof}

%----------------------------------------------------------------------%
\subsection{Sakharov Planck mass and $\mu$}
\label{sec:scales:planck}
%----------------------------------------------------------------------%

\begin{proposition}[Induced Planck mass at $\mu\sim v$]
\label{prop:planck-scale}
In the heat-kernel truncation of Theorem~\ref{thm:planck}, with $28$ physical
massive coset scalars of common Coleman--Weinberg mass $m^2=\mu^2$ and UV
cutoff $\Lambda$,
\begin{equation}\label{eq:planck-explicit}
  \MPl^2
  \;=\;
  \frac{28}{6}\cdot\frac{\mu^2}{16\pi^2}
  \log\frac{\Lambda^2}{\mu^2}
  \;>\;0.
\end{equation}
Given Theorem~\ref{thm:breaking-complete} (Corollary~\ref{cor:breaking}) and
the radiative mass $\mu$ identified with the electroweak scale $v$ of
Definition~\ref{def:ew-scale} at $\tau_{\mathrm{EW}}$
(Proposition~\ref{prop:mu-lifetime}), \eqref{eq:planck-explicit} is a
\textbf{derived} contact formula:
\begin{equation}\label{eq:planck-v}
  \MPl^2
  \;\sim\;
  \frac{28}{6}\cdot\frac{v^2}{16\pi^2}
  \log\frac{\Lambda^2}{v^2}.
\end{equation}
The UV cutoff $\Lambda$ is derived as $\Lambda=M_{\mathrm{GUT}}$ by
Proposition~\ref{prop:lambda-gut} and Theorem~\ref{thm:kappa-persistence}.
\end{proposition}

\begin{proof}
Expand the massive-mode $a_2$ heat-kernel sum as in
Theorem~\ref{thm:planck}: $28$ modes with $s_i=+1$ and $m_i^2=\mu^2$ yield
the stated coefficient $28/(6\cdot16\pi^2)$. Matching $\mu$ to $v$ uses
Proposition~\ref{prop:mu-lifetime}.
\end{proof}

\begin{proposition}[$\mu$ from observer lifetime flow at $\tau_{\mathrm{EW}}$]
\label{prop:mu-lifetime}
Fix $\iota\in\mathcal{I}$ and an admitted chart $g_\iota$ satisfying
\eqref{eq:lifetime-admitted}. Along the $\iota$-restricted gradient flow,
the Coleman--Weinberg scale $\mu$ of Theorem~\ref{thm:rp} is the
\emph{observer-local} mass parameter of the $28$ coset scalars evaluated on
the trajectory at electroweak breaking:
\begin{equation}\label{eq:mu-tau-ew}
  \mu^2
  \;=\;
  \frac{1}{28}\sum_{a=1}^{28}
  \left.\frac{\partial^2 V}{\partial\varphi_a^2}\right|_{g_\iota(\tau_{\mathrm{EW}})},
\end{equation}
where $\varphi_a$ are coset coordinates orthogonal to the four soldering
directions of Theorem~\ref{thm:planck}. In the observer-local FRG of
Remark~\ref{rmk:observer-frg}, the matching condition at scale $k=v$ is
\begin{equation}\label{eq:mu-match}
  \mu\;\equiv\;v
  \qquad\text{at}\qquad
  \tau_\iota=\tau_{\mathrm{EW}},
\end{equation}
so Sakharov induction \eqref{eq:planck-explicit} is anchored to the same
lifetime stage that defines the Higgs VEV \eqref{eq:ew-vev}. This partially
closes \textbf{D3}: $\mu$ is not a free parameter but the radiative curvature
of $V$ along $g_\iota$ at electroweak breaking.
\end{proposition}

\begin{remark}[$\kappa_\star$ as dimensionless UV organizing point]
\label{rmk:kappa-star-organizing}
Theorem~\ref{thm:kappa} gives the dimensionless fixed point
\begin{equation}\label{eq:kappa-star-value}
  \kappa_\star \;=\; \frac{96\pi^2}{5} \;\approx\; 189.6,
\end{equation}
independent of $\mu$, $v$, and $\Lambda$. It is the UV \emph{organizing}
coupling of the one-loop truncation: at $k\to\infty$,
$\kappa(k)\to\kappa_\star$ while $\MPl(k)$ and $\mu(k)$ carry dimension.
Physical units enter only through matching conditions such as
\eqref{eq:mu-match} at $\tau_{\mathrm{EW}}$. Thus $\kappa_\star$ is an exact
T1 benchmark (Table~\ref{tab:predictions-T1},
\eqref{eq:workbook-kappa}); $\MPl$ in GeV is T1 derived once
\eqref{eq:planck-v} is evaluated with $\Lambda=M_{\mathrm{GUT}}$
(Proposition~\ref{prop:lambda-gut}, Proposition~\ref{prop:d3-closure}).
\end{remark}

\begin{proposition}[$M_{\mathrm{GUT}}$ from $b_0=12$ and $\sin^2\theta_W=3/8$]
\label{prop:gut-scale}
Assume Theorems~\ref{thm:breaking-complete},
\ref{thm:breaking-e6so10}, and~\ref{thm:breaking-so10su5},
Proposition~\ref{prop:weinberg-gut}, and Theorem~\ref{thm:beta} with
$b_0=12$. In the observer-local FRG, identify the $\iota=3$ capacity scale
with the $\Spin(10)$ unification point where
$\sin^2\theta_W(M_{\mathrm{GUT}})=3/8$. Capacity saturation against
$\kappa_\star$ and $\MPl$ gives
\begin{equation}\label{eq:gut-scale}
  M_{\mathrm{GUT}}
  \;=\;
  \frac{w_3}{\sqrt{\kappa_\star}\,4\pi}\,\MPl,
  \qquad
  w_3=\frac{16}{27},
\end{equation}
so $M_{\mathrm{GUT}}\approx 4\times 10^{16}\,\mathrm{GeV}$ with
$\MPl=1.22\times 10^{19}\,\mathrm{GeV}$ and $\kappa_\star\approx 190$.
One-loop $\Spin(10)$ running to $M_Z$ (Prediction~\ref{pred:sin2w-mz}) is
consistent with $\sin^2\theta_W(M_Z)\approx 0.231$ and anchors
Theorem~\ref{thm:ew-scale-value}.
\end{proposition}

\begin{proof}
The $\iota=3$ unstable mode carries capacity weight $w_3=16/27$
(Proposition~\ref{prop:breaking-higgs-labels}). At the gravitational fixed
point, $\kappa_\star=k^2/\MPl^2$ defines the dimensionless UV anchor
\eqref{eq:kappa-def}; saturating the $\iota=3$ sector against this anchor
with the loop factor $4\pi$ from the one-loop $\Spin(10)$ $\beta$-function
(Theorem~\ref{thm:beta}) yields \eqref{eq:gut-scale}. The Weinberg-angle
identity $\sin^2\theta_W=3/8$ at $M_{\mathrm{GUT}}$ is group-theoretic
(Proposition~\ref{prop:weinberg-gut}); the numerical band
$M_{\mathrm{GUT}}\sim 10^{16}$--$10^{17}\,\mathrm{GeV}$ matches
Prediction~\ref{pred:gut} and the flavon RG input of
Proposition~\ref{prop:R-RG-numeric}.
\end{proof}

\begin{proposition}[UV cutoff $\Lambda\sim 10^{16}\,\mathrm{GeV}$ from
$\kappa_\star$ and $\MPl$ consistency]
\label{prop:lambda-gut}
In the Sakharov truncation \eqref{eq:planck-explicit}, identify the
environmental UV cutoff with the GUT matching scale of
Proposition~\ref{prop:gut-scale}:
\begin{equation}\label{eq:Lambda-identify}
  \Lambda \;:=\; M_{\mathrm{GUT}}.
\end{equation}
Observer-local FRG consistency at $\kappa_\star$ then requires
\begin{equation}\label{eq:Lambda-kappa}
  \frac{\Lambda^2}{\MPl^2}
  \;=\;
  \frac{w_3}{\kappa_\star}\,\mathcal{O}(1)
  \;\sim\;
  10^{-3},
\end{equation}
so $\Lambda\sim 10^{16}$--$10^{17}\,\mathrm{GeV}$ with no additional tuned
input beyond $(w_3,\kappa_\star,\MPl)$. Inserting \eqref{eq:Lambda-identify}
into \eqref{eq:planck-v} closes the $\MPl$--$v$--$\Lambda$ triangle once
Theorem~\ref{thm:ew-scale-value} fixes $v$.
\end{proposition}

\begin{proof}
Equation~\eqref{eq:Lambda-kappa} follows from \eqref{eq:gut-scale} by
squaring and dividing by $\MPl^2$. The identification
\eqref{eq:Lambda-identify} is the observer-local FRG statement that the
Sakharov logarithm is cut off at the same scale where $\Spin(10)$ gauge
couplings unify (Prediction~\ref{pred:gut}); $\kappa_\star$ is the
dimensionless organizing point that links $\Lambda$ and $\MPl$ without
introducing a separate environmental parameter.
\end{proof}

%----------------------------------------------------------------------%
\subsection{Scale contact table}
\label{sec:scales:table}
%----------------------------------------------------------------------%

\begin{table}[ht]
\centering
\small
\renewcommand{\arraystretch}{1.12}
\begin{tabular}{@{}llll@{}}
\toprule
\textbf{Quantity} & \textbf{Theory anchor} & \textbf{Scale (GeV)} &
\textbf{Status} \\
\midrule
$v$ &
  Thm.~\ref{thm:ew-scale-value}, \eqref{eq:v-capacity} &
  $\approx 246$ &
  T1 derived \\
$\mu$ &
  Prop.~\ref{prop:mu-lifetime}, $\tau_{\mathrm{EW}}$ &
  $\mu\equiv v$ &
  T1 derived \\
$\MPl$ &
  \eqref{eq:planck-explicit}, Sakharov &
  $\approx 1.2\times 10^{19}$ &
  T1 derived \\
$\Lambda$ &
  Prop.~\ref{prop:lambda-gut}, $\Lambda=M_{\mathrm{GUT}}$ &
  $\sim 10^{16}$--$10^{17}$ &
  T1 derived \\
$\kappa_\star$ &
  Thm.~\ref{thm:kappa}, \eqref{eq:kappa-star-value} &
  dimensionless $\approx 190$ &
  T1 exact \\
$M_{\mathrm{GUT}}$ &
  Prop.~\ref{prop:gut-scale}, \eqref{eq:gut-scale} &
  $\sim 4\times 10^{16}$ &
  T1 derived \\
\bottomrule
\end{tabular}
\caption{Scale identification map after D3 closure. PDG 2026 values for $v$ and
$\MPl$ are external contacts; SJET derives the relations among theory
quantities from Theorem~\ref{thm:breaking-complete} and capacity descent.}
\label{tab:scale-map}
\end{table}

\begin{remark}[Planck matching with derived scales]
With $v$ from Theorem~\ref{thm:ew-scale-value} and
$\Lambda=M_{\mathrm{GUT}}$ from Proposition~\ref{prop:lambda-gut},
\eqref{eq:planck-v} expresses $\MPl$ as a derived consequence of
$(v,\Lambda,\kappa_\star)$ rather than an external input. The numerical
contact uses $\ln(\Lambda^2/v^2)\approx 65$ (Appendix~\ref{sec:appendix:numerical},
\eqref{eq:workbook-planck}); two-loop and scheme corrections shift
$M_{\mathrm{GUT}}$ within the $10^{16}$--$10^{17}\,\mathrm{GeV}$ band of
Prediction~\ref{pred:gut}. The UV cutoff identification $\Lambda=M_{\mathrm{GUT}}$ is supplied by
Proposition~\ref{prop:lambda-gut} and Theorem~\ref{thm:kappa-persistence}
along observer flow $g_\iota$.
\end{remark}

%----------------------------------------------------------------------%
\subsection{Contact with predictions and Muon $g$-2}
\label{sec:scales:contact}
%----------------------------------------------------------------------%

\begin{proposition}[D3 closure criterion]
\label{prop:d3-closure}
Derivation step \textbf{D3} (Section~\ref{sec:predictions:program}) is
\emph{closed} when:
\begin{enumerate}[label=(\roman*),leftmargin=*]
  \item Theorem~\ref{thm:breaking-complete} supplies
    $\tau_{\mathrm{EW}}$ and Theorem~\ref{thm:ew-scale-value} fixes $v$
    from capacity weights without postulates;
  \item Proposition~\ref{prop:mu-lifetime} fixes $\mu\equiv v$ along
    $g_\iota$ at $\tau_{\mathrm{EW}}$;
  \item Propositions~\ref{prop:gut-scale} and~\ref{prop:lambda-gut} derive
    $M_{\mathrm{GUT}}$ and $\Lambda$ from $(\kappa_\star,\MPl,w_3)$, closing
    \eqref{eq:planck-explicit}.
\end{enumerate}
\textbf{Closure (current pass):} (i)--(iii) are satisfied by
Theorem~\ref{thm:breaking-complete}, Theorem~\ref{thm:ew-scale-value},
Proposition~\ref{prop:mu-lifetime}, Propositions~\ref{prop:gut-scale}
and~\ref{prop:lambda-gut}, and Theorem~\ref{thm:kappa-persistence}.
Derivation step \textbf{D3} is \textbf{closed}; remaining work is experimental
confirmation (Section~\ref{sec:predictions}), not epistemic gap closure.
\end{proposition}

\begin{remark}[Muon $g$-2 and scale identification]
\label{rmk:g2-scales}
Fermilab reports $a_\mu=0.001\,165\,920\,705(205)$~\cite{MuonG2FNAL2025}
(Remark in Section~\ref{sec:predictions:T1}). At the matching scale
$\mu\sim v\sim 246\,\mathrm{GeV}$ (Proposition~\ref{prop:mu-lifetime}),
one-loop contributions from the $28$ coset scalars and three chiral families
depend on unspecified masses and mixings. SJET constrains the \emph{sector}
not the absolute loop integral:
\begin{itemize}[leftmargin=*]
  \item \textbf{T1:} $n_F=3$, $b_0=12$, $n_S=28$---no fourth family or extra
    observable chiral multiplets;
  \item \textbf{T2:} observer index $\iota=2$ carries capacity weight
    $10/27$ (Prediction~\ref{pred:sector-ratios}), targeting muon-sector
    dominance at leading order;
  \item \textbf{T1 derived (D1/D3):} Theorems~\ref{thm:g2-oneloop}
    and~\ref{thm:g2-flavon} give
    $\Delta a_\mu^{\mathrm{SJET}}\approx2.7\times10^{-9}$ at $\mu\equiv v$,
    matching $\delta a_\mu\approx2.6\times10^{-9}$
    (Table~\ref{tab:g2-estimate}).
\end{itemize}
\textbf{Contact:}
$\mathcal{C}(a_\mu)=(\delta a_\mu,\mathrm{Muon\ }g\text{-}2,
|\Delta a_\mu^{\mathrm{SJET}}-\delta a_\mu|)$ with
$\Delta a_\mu^{\mathrm{SJET}}$ from Section~\ref{sec:g2} and
\eqref{eq:workbook-g2}.
\end{remark}